\section{Identify and materialize the instance}\label{sec:materialization}

A static checker consumes a matrix $c$ of costs from vertex $i$ to vertex $j$, whereas a TD checker consumes a common departure horizon and, for each directed arc, an ordered sequence of knots of its ATF. We call the process of obtaining these values from the published instance \emph{materialization}. Once the binary64 values are fixed, a checker reading an entire JSON sidecar and a checker regenerating values on demand face the same evaluation problem.

\paragraph{Identify the published inputs.}
The \emph{instance manifest} is the record that tells a consumer which files and numerical rules belong to this edition. It \must{} identify the family and edition, problem type, objective, units, default profile, supported features, and the descriptor and required sidecars by path and SHA-256. The descriptor supplies the depot and customer indices, demands, capacity, fleet limit and applicable temporal data. A result refers to its manifest by the SHA-256 of its file bytes. Appendix~\ref{app:certification-order} gives the publication order and explains how to verify an edition outside a Git checkout. The customer count determines the expected matrix dimensions or directed arc set. In this core, demands and vehicle capacity are nonnegative integers exactly representable in binary64, and the depot has zero demand and zero service duration. A supplied fleet limit is a nonnegative integer, and an absent limit imposes no bound. Since changing a time window can change feasibility while preserving all arc values, the result identifies the complete manifest.

For a static instance, \id{units.cost} names the unit of the arc matrix. In a VRPTW the same matrix also supplies the travel times, so \id{units.time} uses the same unit for windows and service durations. In a CVRP the checker uses the matrix only for the objective. For example, the Mamut2026 fastest and shortest CVRP variants measure travel time in seconds and distance in metres, respectively. Both have the same capacity constraints.

A family descriptor resolves its relative sidecar paths from the family collection root, which is marked by \id{mamut-collection.json}. For example, the Rifki-30 descriptor in \path{TDVRP/n=30/} refers to \path{sidecars/n=30/Rifki-30.atf.json.gz} under the Rifki2020 root. A standalone example carrying its own sidecar instead resolves the path from its descriptor's directory. For a compressed artifact, the manifest \must{} specify whether a digest covers the distributed compressed bytes or the decompressed content, because recompressing identical content can change the former digest. The published \href{https://github.com/ANR-MAMUT/MAMUT-routing-lib/blob/e032f2b9b9d1ba16a561325a07ab59cf39f8e5d5/src/mamut_routing_lib/sidecars.py}{collection resolver} documents the path rules.

\paragraph{Choose the source of arc values.}
The manifest's \id{materialization} object identifies the representation and its \id{interpretation}, here \id{binary64}. A static instance can use an \id{embedded-matrix} or a \id{sidecar-matrix}, and a TD instance can supply a \id{stored-atf} sidecar. In these cases, the checker reads the published entries or knots directly. A family using \id{regenerated} data \must{} identify the raw inputs by content hash, the materialization function and immutable implementation, the numerical parameters and the relevant execution environment. Record the callable, identify its preserved source distribution by content hash and pin the dependencies that affect numerical output, including parameters for rounding, path selection or function simplification where used. A versioned environment record may be referenced by hash across many manifests.

For example, a Mamut2026 base can yield both a shortest-distance matrix and a fastest-travel-time matrix, and path selection as well as the subsequent 3-decimal rounding are part of the respective materialization recipes. For a TD road model, the traffic overlay can change the arrival function on a path even when the road graph and vertices stay fixed. Other TD families derive their ATFs from speed periods and arc categories, or distribute already computed knots. In each case, the fixed operands reaching the checker are the output of the declared recipe. In particular, coordinates included for display do not authorize replacing road costs by Euclidean distances \autocite{pichonMamut2026,pichonPoryos2026}.

\paragraph{Verify the complete materialization.}
A raw SHA-256 checks the identity of a particular file, whereas a \emph{semantic digest} checks the complete sequence of interpreted binary64 operands. For a static matrix, the manifest field is \id{materialization.matrix_sha256}, which covers the dimension and all entries in row-major order, including the diagonal. For TD data, the field is \id{materialization.atf_stream_sha256}, which covers the vertex count, horizon and ordered knots on every directed arc between distinct vertices. Preserving the knot order matters at a jump, where the first and last ordinates at one time have different roles. In released TD descriptors, the separate field \id{td.atf_sha256} covers the canonical ATF JSON content, including its metadata. Appendix~\ref{app:materialization-format} gives the exact byte streams for the two semantic digests and companion examples.

A verifier \must{} compare the complete interpreted stream with the declared semantic digest, although it may produce rows or arcs incrementally and discard each after hashing. A checker may reuse a successful full verification while the raw inputs, code, environment, parameters, interpretation and expected digest remain unchanged. A checker evaluating one submitted solution can then request just its traversed arcs, provided that the on-demand call produces the corresponding entries of the verified materialization. Note that hashing that subset alone only establishes which operands the check used, whereas checking the identity of the whole instance requires the complete stream. A missing required input, digest mismatch or unavailable required materializer stops certification.

Historical generation documentation records how the publisher obtained the descriptor, raw inputs or stored sidecars, including the generator revision and parameters when known. An imported dataset may have a precisely identified source file even when part of its earlier construction is unknown. By contrast, the information needed to read or regenerate the authoritative arc values belongs to the instance contract, because every checker must be able to recover those values.
